\section{Zadanie 15}
\autor{Filip Manijak}
Udowodnij, że każdy graf dwudzielny z $m$ krawędziami zawiera podgraf na
co najmniej $\frac{3}{4} m^{\frac{2}{3}}$ krawędziach, który nie zawiera cyklu na czterech
wierzchołkach $C_4$ jako podgrafu.

\vspace{0.5em}

\textbf{Uwaga:} Wystarczy pokazać ograniczenie $\frac{1}{2} m^{\frac{2}{3}}$.

\vspace{0.5em}

\textbf{Wskazówka:} Grafy można losować. Jednym ze sposobów jest ustalenie
pewnego parametru $p$ i dla każdej pary wierzchołków niezależne wylosowanie
tego, czy łączy je krawędź, czy nie, z prawdopodobieństwem $p$. Jeśli
losowany graf nie działa, to może należy go zmodyfikować po wylosowaniu?


\subsection*{Rozwiązanie}
Najpierw pokażę lemat:
\subsection{Lemat: Liczba cykli w grafie dwudzielnym to max $\frac{|E|^2}{4}$}
Dlaczego? Bo par rozłącznych krawędzi jest $\frac{|E||E-1|}{2}$. Zauważmy, że każda para rozłącznych krawędzi występuje razem w max. 1 $C_4$. Oszacowanie jeszcze musimy podzielić na 2, bo każdy cykl zliczamy 2 razy. 
\subsection{Główny pomysł}
Będziemy losować graf (każdą krawędź z pwdp. $p$) i zauważymy że liczba cykli jest całkiem mała. W podgrafie naszego grafu dwudzielnego, powstałym w taki sposób że każda krawędź występuje w nim z prawdopodobieństwem $p$ jest w wartości oczekiwanej:
\begin{itemize}
    \item $pm$ krawędzi (liniowość wartości oczekiwanej)
    \item $p^4num\_C_4$ $C_4$ (wszystkie 4 krawędzie muszą zostać przeniesione)
\end{itemize}
Zauważmy, że wystarczy usunąć tyle krawędzi co cykli, żeby żadnych cykli nie było. Czyli, stwórzmy nowy sposób losowania grafu: Losujemy każdą krawędź z prawdopodobieństwem $p$, ale jeśli krawędź ta tworzy $C_4$ to jej nie bierzemy. Teraz policzmy wartość oczekiwaną:
$$E[E(H)] = pm - p^4\frac{m^2}{4}$$
Co po podstawieniu $p = m^{-1/3}$ daje:
$$E[E(H)] = \frac{3}{4} m^\frac{2}{3}$$
Z trywialnej obserwacji$^{TM}$ koniec dowodu.

